\section{The irreducible common-mode blind spot: large but finite}
\label{sec:blind}

Suppose an independent anchor is present, so the common-mode statistic of
Section~\ref{sec:gls} is live and the pre-fit residual carries a fault along some
unit direction $d\in\R^N$, of magnitude $\xi$. The honest question is: what is the
largest $\xi$ that \emph{neither} the whitened separation detector \emph{nor} the
common-mode detector can see? This section answers it as a \Proven{} bound and
states the blind spot plainly. The answer (caveat~I1) is the paper's honesty
centerpiece: the blind spot is large but \emph{finite}, and our ceiling
under-states it.

\subsection{Two detectors, two blind directions}
\label{sec:twodetectors}

Write the undetectable ceiling along $d$ as the smaller of the two magnitudes at
which the two detectors respectively trip,
\begin{equation}
  \xi^\star(d) \;=\; \min\big(\alpha_{\mathrm{ss}}(d),\ \alpha_{\mathrm{cm}}(d)\big),
  \label{eq:ceiling}
\end{equation}
where $\alpha_{\mathrm{ss}}$ is the trip magnitude of an idealized whitened parity
(solution-separation) detector and $\alpha_{\mathrm{cm}}$ that of the common-mode
statistic. A detector \emph{blind} to $d$ has an infinite trip magnitude along it.
The two blind conditions are, from the parity projection and
from~\eqref{eq:cmstat} respectively (Appendix~\ref{app:blind}),
\begin{equation}
  \alpha_{\mathrm{ss}}(d)=\infty \iff d \parallel \bmone{},
  \qquad
  \alpha_{\mathrm{cm}}(d)=\infty \iff \bmone^{\!\top}\Om^{-1} d = 0 .
  \label{eq:blindconds}
\end{equation}
Separation is blind exactly to the common direction $d\parallel\bmone{}$;
common mode is blind exactly to directions $\Om^{-1}$-orthogonal to $\bmone{}$.

\subsection{The blind conditions are mutually exclusive under a PD whitening}
\label{sec:mutex}

\begin{theorem}[Finite ceiling for every nonzero direction]
\label{thm:finite}
If $\Om{}\succ0$ (positive definite), the two blind conditions
of~\eqref{eq:blindconds} cannot both hold for any nonzero $d$; hence the two
detectors jointly cover every nonzero direction and the ceiling~\eqref{eq:ceiling}
is finite for every nonzero $d$. The ceiling is infinite only for degenerate
inputs---$d=0$, or a non-PD (rank-deficient) $\Om{}$.
\end{theorem}

\begin{proof}[Proof sketch]
If $d\parallel\bmone{}$, say $d=\lambda\bmone{}$ with $\lambda\neq0$, then
$\bmone^{\!\top}\Om^{-1}d=\lambda\,\bmone^{\!\top}\Om^{-1}\bmone{}>0$ because
$\Om^{-1}\succ0$ forces the quadratic form $\bmone^{\!\top}\Om^{-1}\bmone{}>0$;
so $\alpha_{\mathrm{cm}}$ is finite. Conversely if
$\bmone^{\!\top}\Om^{-1}d=0$ then $d$ cannot be a nonzero multiple of $\bmone{}$
(such a multiple gives the strictly positive form above); so $\alpha_{\mathrm{ss}}$
is finite. The two blind sets therefore intersect only at $d=0$, and their
complement union is all of $\R^N\setminus\{0\}$. Full detail in
Appendix~\ref{app:blind}.
\end{proof}

\begin{remark}[\textbf{Caveat I1: the blind spot is large but finite, and the
ceiling is a lower bound.}]
\label{rem:i1}
Theorem~\ref{thm:finite} is the referee-facing honesty statement, and it cuts in
the conservative direction. \textbf{The irreducible common-mode blind spot is
large but \emph{finite}: under a positive-definite whitening the separation and
common-mode blind directions are mutually exclusive, so the two detectors jointly
cover every nonzero direction, and the undetectable ceiling is finite for every
nonzero $d$ (infinite only for the degenerate $d=0$ or a non-PD $\Om{}$).} What
remains is not a total or infinite blind spot but a bounded one, driven by
\emph{model mismatch}: the residual undetectable magnitude is largest where the
true fault-coupling direction is nearly aligned with a blind direction of the
\emph{modelled} $\Om{}$, and it grows without bound only as $\Om{}$ is pushed toward
singularity. Two further honesty points sharpen the bound rather than soften it.
First, the $\alpha_{\mathrm{ss}}$ used in~\eqref{eq:ceiling} is an \emph{idealized}
single whitened-contrast parity detector, which is \emph{more} sensitive than a
real integrity-risk-allocated MHSS subset detector with a finite false-alert
budget; a real monitor trips later, so the true undetectable magnitude is at least
$\xi^\star(d)$ and our ceiling is a conservative \emph{lower bound}, never an
overstatement of detectability. Second, we make no claim that the blind spot is
small---only that it is finite and bounded by the model-mismatch between the true
coupling and the modelled $\Om{}$. We never assert an infinite or unqualified blind
spot, and we never assert its absence.
\end{remark}
